\documentclass[aspectratio=169,11pt]{beamer}

\usepackage[T1]{fontenc}
\usepackage[utf8]{inputenc}
\usepackage{lmodern}
\usepackage{microtype}
\usepackage{amsmath,amssymb,mathtools,bm}
\usepackage{booktabs,array}
\usepackage{xcolor}

\definecolor{TitleBlue}{HTML}{163A5F}
\definecolor{AccentTeal}{HTML}{2D6F73}
\definecolor{DeepNavy}{HTML}{102A43}
\definecolor{SoftBlue}{HTML}{EAF3F8}
\definecolor{SoftTeal}{HTML}{EDF8F6}
\definecolor{SoftGray}{HTML}{F4F6F8}
\definecolor{BodyGray}{HTML}{243B53}
\definecolor{WarningRed}{HTML}{9B2C2C}

\setbeamercolor{normal text}{fg=BodyGray,bg=white}
\setbeamercolor{frametitle}{fg=TitleBlue,bg=white}
\setbeamercolor{structure}{fg=AccentTeal}
\setbeamercolor{block title}{fg=white,bg=TitleBlue}
\setbeamercolor{block body}{fg=BodyGray,bg=SoftBlue}
\setbeamercolor{block title alerted}{fg=white,bg=WarningRed}
\setbeamercolor{block body alerted}{fg=BodyGray,bg=SoftGray}
\setbeamercolor{block title example}{fg=white,bg=AccentTeal}
\setbeamercolor{block body example}{fg=BodyGray,bg=SoftTeal}

% User requested title size comparable to the body size.
\setbeamerfont{normal text}{size=\fontsize{12.2}{14.6}\selectfont}
\setbeamerfont{frametitle}{size=\fontsize{12.2}{14.6}\selectfont,series=\bfseries}
\setbeamerfont{block title}{size=\fontsize{11.4}{13.5}\selectfont,series=\bfseries}
\setbeamerfont{block body}{size=\fontsize{10.9}{13.0}\selectfont}
\setbeamerfont{footline}{size=\fontsize{8.6}{10.0}\selectfont}

\setbeamertemplate{navigation symbols}{}
\setbeamertemplate{itemize item}{\color{AccentTeal}$\blacktriangleright$}
\setbeamertemplate{itemize subitem}{\color{AccentTeal}$\bullet$}
\setlength{\leftmargini}{1.05em}
\setlength{\leftmarginii}{1.0em}
\setlength{\parskip}{0.03em}

\newcommand{\bracket}[2]{\left[#1,#2\right]}
\newcommand{\dd}{\,\mathrm{d}}
\newcommand{\avg}[1]{\left\langle #1\right\rangle}
\newcommand{\Rpsi}{R_{\psi}}
\newcommand{\Romega}{R_{\omega}}
\newcommand{\vect}[1]{\bm{#1}}
\newcommand{\code}[1]{\texttt{#1}}
\newcommand{\OmegaD}{\Omega=[0,2\pi)\times[0,2\pi)}
\newcommand{\stateop}{(\psi_0,U_0,\eta,\nu,t)\mapsto(\psi_t,U_t)}
\newcommand{\currentmodule}{Module 1}
\newcommand{\moduleprogress}{1 of 4}
\newcommand{\setmodule}[2]{%
  \renewcommand{\currentmodule}{#1}%
  \renewcommand{\moduleprogress}{#2}%
}

\setbeamertemplate{frametitle}{%
  \vspace{0.06cm}
  \usebeamerfont{frametitle}\insertframetitle\par
  \vspace{-0.20cm}
  {\color{AccentTeal}\rule{\textwidth}{0.55pt}}
}

\setbeamertemplate{footline}{%
  \leavevmode
  \hbox{%
    \begin{beamercolorbox}[wd=.48\paperwidth,ht=2.6ex,dp=1.05ex,leftskip=.45cm]{author in head/foot}%
      \usebeamerfont{footline}\color{TitleBlue}\textbf{PINN\_M3DC1}\hspace{0.45em}|\hspace{0.45em}\currentmodule
    \end{beamercolorbox}%
    \begin{beamercolorbox}[wd=.52\paperwidth,ht=2.6ex,dp=1.05ex,rightskip=.45cm plus1fil]{date in head/foot}%
      \usebeamerfont{footline}\color{BodyGray}\moduleprogress\hspace{0.8em}\textbullet\hspace{0.8em}Slide \insertframenumber\ of \inserttotalframenumber
    \end{beamercolorbox}%
  }
}

\setbeamersize{text margin left=0.60cm,text margin right=0.60cm}
\AtBeginDocument{%
  \setlength{\abovedisplayskip}{0.22em}%
  \setlength{\belowdisplayskip}{0.22em}%
  \setlength{\abovedisplayshortskip}{0.16em}%
  \setlength{\belowdisplayshortskip}{0.16em}%
}

\begin{document}

% --------------------------------------------------------------------------
% Module 1: four slides
% --------------------------------------------------------------------------

\setmodule{Module 1}{1 of 4}
\begin{frame}[t]{Module 1: reduced resistive-MHD state}
\vspace{-0.15em}
The project physics is a strict two-dimensional, incompressible, doubly periodic reduced-MHD model:
\[
\OmegaD,\qquad t\in[0,T],\qquad \partial_z=0.
\]
The primary scalar fields are magnetic flux and stream function,
\[
\vect q(x,y,t)=(\psi,U),
\]
with derived quantities
\[
\vect B_\perp=\nabla\psi\times\widehat{\vect z}=(\psi_y,-\psi_x),\qquad
\vect v_\perp=\widehat{\vect z}\times\nabla U=(-U_y,U_x),
\]
\[
J=-\nabla^2\psi,\qquad \omega=\nabla^2U,
\qquad \nabla\cdot\vect B_\perp=\nabla\cdot\vect v_\perp=0.
\]
The revised learning direction uses the same physics, but the surrogate input becomes a complete initial field pair rather than only coordinates and scalar case parameters.
\end{frame}

\setmodule{Module 1}{2 of 4}
\begin{frame}[t]{Module 1: brackets and governing equations}
\vspace{-0.15em}
The two-dimensional bracket is
\[
\bracket{f}{g}=f_xg_y-f_yg_x,
\qquad
\vect v_\perp\cdot\nabla f=\bracket{U}{f}.
\]
Thus nonlinear advection appears compactly as a Poisson bracket. The frozen sign convention is
\[
\boxed{\psi_t+\bracket{U}{\psi}=\eta\nabla^2\psi,}
\]
\[
\boxed{\omega_t+\bracket{U}{\omega}=\bracket{J}{\psi}+\nu\nabla^2\omega,}
\qquad
J=-\nabla^2\psi,
\quad
\omega=\nabla^2U.
\]
Interpretation:
\[
\bracket{U}{\psi}:\text{ flux advection},\quad
\eta\nabla^2\psi=-\eta J:\text{ resistive slippage},
\]
\[
\bracket{J}{\psi}:\text{ curl of Lorentz force},\quad
\nu\nabla^2\omega:\text{ viscous vorticity diffusion}.
\]
\end{frame}

\setmodule{Module 1}{3 of 4}
\begin{frame}[t]{Module 1: scales, gauges, and energy}
\vspace{-0.15em}
Alfven normalization gives
\[
V_A=\frac{B_0}{\sqrt{\mu_0\rho_0}},
\qquad
 t_A=\frac{L_0}{V_A},
\qquad
\eta=\frac{\eta_m}{L_0V_A},
\qquad
\nu=\frac{\nu_d}{L_0V_A}.
\]
Dimensionless control numbers are
\[
S=\eta^{-1},\qquad \mathrm{Re}=\nu^{-1},\qquad \mathrm{Pm}=\frac{\nu}{\eta}.
\]
Periodic gauges and means:
\[
\avg{U}(t)=0,
\qquad
\frac{\mathrm d}{\mathrm dt}\avg{\psi}=0,
\qquad
\avg J=\avg\omega=0.
\]
The verification energy law is
\[
\boxed{\frac{\mathrm d}{\mathrm dt}\frac12\int_\Omega(|\nabla\psi|^2+|\nabla U|^2)\dd A
=-\eta\int_\Omega J^2\dd A-\nu\int_\Omega\omega^2\dd A\le0.}
\]
\end{frame}

\setmodule{Module 1}{4 of 4}
\begin{frame}[t]{Module 1: claim boundary under the operator change}
\vspace{-0.15em}
The physical model does \emph{not} change when the neural architecture changes. What changes is the learned map:
\[
\text{old coordinate surrogate:}\qquad (x,y,t,\eta,\nu,A_0)\mapsto(\psi,U),
\]
\[
\text{new full-field operator:}\qquad
\boxed{\stateop.}
\]
In Module 1 language, this means:
\[
(\psi_0,U_0)\text{ carries the initial topology and flow shape},
\qquad
(\eta,\nu)\text{ carry transport strength}.
\]
\[
A_0\text{ may remain metadata, but in version 1 it is already encoded through }U_0.
\]
The scientific boundary remains strict: these notes support a synthetic two-field reduced-MHD surrogate workflow, not a validated surrogate of genuine M3D-C1 output.
\end{frame}

% --------------------------------------------------------------------------
% Module 2: four slides
% --------------------------------------------------------------------------

\setmodule{Module 2}{1 of 4}
\begin{frame}[t]{Module 2: solver output becomes operator data}
\vspace{-0.15em}
For each physical case
\[
\vect\mu=(\eta,\nu,A_0),
\]
the verified solver produces a full trajectory
\[
\mathcal Q_\mu=\{\psi(x,y,t),U(x,y,t),J(x,y,t),\omega(x,y,t): (x,y)\in\Omega,
0\le t\le T\}.
\]
The integrated revision makes the downstream handoff explicit:
\[
\boxed{X_k=(\psi[0,:,:],U[0,:,:],\eta,\nu,t_k),
\qquad
Y_k=(\psi[k,:,:],U[k,:,:]).}
\]
These are raw numerical scalar arrays. Contours, field-line plots, PNGs, screenshots, and rendered images are excluded from the learning-data path.
\end{frame}

\setmodule{Module 2}{2 of 4}
\begin{frame}[t]{Module 2: island coalescence initial condition}
\vspace{-0.15em}
The benchmark starts from
\[
\psi_0(x,y)=\Psi_{\rm eq}\sin x\sin y,
\qquad
\Psi_{\rm eq}=1,
\]
\[
\boxed{U_0(x,y)=A_0(\cos x-\cos y).}
\]
Derived initial fields are
\[
J_0=-\nabla^2\psi_0=2\psi_0,
\qquad
\omega_0=\nabla^2U_0=-U_0,
\]
\[
\vect B_{\perp,0}=(\sin x\cos y,-\cos x\sin y),
\qquad
\vect v_{\perp,0}=(-A_0\sin y,-A_0\sin x).
\]
Important limitation for operator learning:
\[
U_0^{(A_0)}(x,y)=A_0 f(x,y)
\quad\Rightarrow\quad
\text{one fixed initial-flow shape with varying amplitude.}
\]
\end{frame}

\setmodule{Module 2}{3 of 4}
\begin{frame}[t]{Module 2: pseudo-spectral reference calculation}
\vspace{-0.15em}
The solver advances the Fourier state
\[
\widehat{\vect q}(t)=\begin{bmatrix}\widehat\psi(t)\\ \widehat\omega(t)\end{bmatrix},
\qquad
\widehat U=-\frac{\widehat\omega}{k^2}\;(k\ne0),
\qquad
\widehat J=k^2\widehat\psi.
\]
Spectral derivatives are algebraic:
\[
\widehat{f_x}=ik_x\widehat f,
\quad
\widehat{f_y}=ik_y\widehat f,
\quad
\widehat{\nabla^2 f}=-k^2\widehat f.
\]
Nonlinear brackets are dealiased physical-grid products:
\[
\widehat{\mathcal B(f,g)}=\mathcal P_{2/3}\mathcal F(f_xg_y-f_yg_x),
\qquad
|m|\le\lfloor N_x/3\rfloor,
\quad
|n|\le\lfloor N_y/3\rfloor.
\]
Classical RK4 advances the stage-consistent state; $(U,J)$ are reconstructed at every stage, not reused from stale fields.
\end{frame}

\setmodule{Module 2}{4 of 4}
\begin{frame}[t]{Module 2: acceptance before pair extraction}
\vspace{-0.15em}
A trajectory is accepted only after verification and numerical adequacy:
\[
\begin{aligned}
\text{operators}&\rightarrow\text{analytic fields}\rightarrow\text{brackets}\rightarrow\text{exact decay}\\
&\rightarrow\text{manufactured solution}\rightarrow\text{refinement and parity}.
\end{aligned}
\]
Physics diagnostics include
\[
E=\frac12\int_\Omega(|\nabla\psi|^2+|\nabla U|^2)\dd A,
\quad
D_\eta=\eta\int_\Omega J^2\dd A,
\quad
D_\nu=\nu\int_\Omega\omega^2\dd A,
\]
\[
J_{\max},\quad \text{sheet width},\quad \text{spectral tails},\quad
\Delta\psi_{OX}=\psi_X-\psi_O,
\quad
\Phi_{\rm rec}=\Delta\psi_{OX}(0)-\Delta\psi_{OX}(t).
\]
One trajectory may yield many target-time operator pairs, but it is still one physical case. Splits must occur before extracting $X_k\mapsto Y_k$ pairs.
\end{frame}

% --------------------------------------------------------------------------
% Module 3: four slides
% --------------------------------------------------------------------------

\setmodule{Module 3}{1 of 4}
\begin{frame}[t]{Module 3: full-field operator example}
\vspace{-0.15em}
Module 3 defines the learning object without modifying canonical solver truth:
\[
\boxed{
\mathcal X_{c,k}=\bigl(\psi_{c,0},U_{c,0},\eta_c,\nu_c,t_k\bigr),
\qquad
\mathcal Y_{c,k}=\bigl(\psi_{c,k},U_{c,k}\bigr).}
\]
Tensor meanings are explicit:
\[
\psi_{c,k},U_{c,k}\in\mathbb R^{N_x\times N_y},
\qquad
q[n,i,j]=q(t_n,x_i,y_j).
\]
Identifiers must remain separate:
\[
\code{physical\_case\_id}\ne\code{case\_id}\ne\code{initial\_state\_id}\ne\code{operator\_example\_id}.
\]
Different times from one trajectory are correlated examples, not independent cases.
\end{frame}

\setmodule{Module 3}{2 of 4}
\begin{frame}[t]{Module 3: HDF5 contract and tensor packing}
\vspace{-0.15em}
Canonical case files store fields before learning views:
\[
\code{fields/psi},\code{fields/U},\code{fields/J},\code{fields/omega}
\in\mathbb R^{N_t\times N_x\times N_y}.
\]
The operator tensor view may pack channels as
\[
\mathcal X_{c,k}\in\mathbb R^{C_{\rm in}\times N_x\times N_y},
\qquad
C_{\rm in}=5:
\quad
[\psi_0,U_0,\eta\mathbf 1,\nu\mathbf 1,t_k\mathbf 1],
\]
\[
\mathcal Y_{c,k}\in\mathbb R^{2\times N_x\times N_y}:
\quad
[\psi_k,U_k].
\]
Broadcast scalar channels are constant fields. No image resizing, color mapping, rasterization, or aspect-ratio distortion is permitted.
\end{frame}

\setmodule{Module 3}{3 of 4}
\begin{frame}[t]{Module 3: splitting, scaling, and leakage}
\vspace{-0.15em}
Splits are made on complete physical trajectories:
\[
\mathcal P_{\rm all}=\mathcal P_{\rm tr}\cup\mathcal P_{\rm va}\cup\mathcal P_{\rm ti}\cup\mathcal P_{\rm te},
\qquad
\mathcal P_a\cap\mathcal P_b=\varnothing.
\]
Only after this step are target times selected:
\[
(c\in\mathcal P_s)\Rightarrow (c,k)\in\mathcal E_s
\quad\text{for all allowed }k.
\]
Training-only transforms are fitted once:
\[
(m_\psi,s_\psi,m_U,s_U,m_\eta,s_\eta,m_\nu,s_\nu)=\operatorname{fit}(\mathcal P_{\rm tr}).
\]
For case-balanced sampling,
\[
\Pr(c)=\frac{1}{|\mathcal P_{\rm tr}|},
\qquad
\Pr(k\mid c)=\text{declared target-time rule}.
\]
Random pair splits would leak nearly identical time-adjacent physics into test sets.
\end{frame}

\setmodule{Module 3}{4 of 4}
\begin{frame}[t]{Module 3: M3D-C1-shaped I/O without false equivalence}
\vspace{-0.15em}
The dataset is M3D-C1-shaped only at the interface level:
\[
\text{coordinates}+	ext{fields}+	ext{parameters}+	ext{diagnostics}+	ext{metadata}+\text{splits}.
\]
It is not genuine M3D-C1 physics. An eventual adapter must declare
\[
\text{geometry},\quad \text{units},\quad \text{mesh},\quad \text{field names},\quad
\text{signs},\quad \text{gauges},\quad \text{closures},\quad \text{projection losses}.
\]
Module 3 release checks include
\[
\begin{array}{lll}
\text{axis parity}, & \text{round trip}, & J+\nabla^2\psi,\\
\omega-\nabla^2U, & \avg U, & \text{split disjointness},\\
\text{no-rendering provenance}. & &
\end{array}
\]
The learning view is reversible bookkeeping; it must not overwrite accepted reference arrays.
\end{frame}

% --------------------------------------------------------------------------
% Module 4: four slides
% --------------------------------------------------------------------------

\setmodule{Module 4}{1 of 4}
\begin{frame}[t]{Module 4: data-only neural-operator baseline}
\vspace{-0.15em}
The baseline is no longer a coordinate MLP. It is a full-field neural operator:
\[
\boxed{\mathcal G_\theta:(\psi_0,U_0,\eta,\nu,t)\mapsto(\widehat\psi_t,\widehat U_t).}
\]
Inputs and outputs are numerical arrays:
\[
\psi_0,U_0,\widehat\psi_t,\widehat U_t\in\mathbb R^{N_x\times N_y}.
\]
A strict data-only baseline answers:
\[
\text{How well can observed trajectory pairs alone train a field-to-field surrogate?}
\]
No PDE residual, initial-condition penalty, boundary loss, energy balance, target $J$, target $\omega$, derivative loss, or spectral loss updates the weights.
\end{frame}

\setmodule{Module 4}{2 of 4}
\begin{frame}[t]{Module 4: Fourier neural operator structure}
\vspace{-0.15em}
A version-1 input tensor can be
\[
\mathcal X=[\psi_0,U_0,\eta\mathbf 1,\nu\mathbf 1,t\mathbf 1]
\in\mathbb R^{5\times N_x\times N_y}.
\]
FNO layers alternate local channel mixing and truncated spectral convolution:
\[
\vect h^{(\ell+1)}(x,y)
=\sigma\!\bigl(W^{(\ell)}\vect h^{(\ell)}(x,y)
+\mathcal F^{-1}[\,R^{(\ell)}(k)\,\mathcal F\vect h^{(\ell)}(k)\,](x,y)\bigr).
\]
The baseline architecture is a scientific comparator, not the final claim:
\[
\text{lift}\rightarrow 4\text{ FNO blocks}\rightarrow \text{projection to }(\psi,U),
\]
with width, retained modes, activation, and depth selected only by validation cases.
\end{frame}

\setmodule{Module 4}{3 of 4}
\begin{frame}[t]{Module 4: strict supervised full-field loss}
\vspace{-0.15em}
For one operator pair $(c,k)$, standardized output errors are fields:
\[
E_{\psi,c,k}=\widehat{\widetilde\psi}_{c,k}-\widetilde\psi_{c,k},
\qquad
E_{U,c,k}=\widehat{\widetilde U}_{c,k}-\widetilde U_{c,k}.
\]
The case-balanced data-only loss is
\[
\boxed{
\mathcal L_{\rm data}=\frac1{|\mathcal C_B|}\sum_{c\in\mathcal C_B}
\frac1{|\mathcal K_c|}\sum_{k\in\mathcal K_c}
\left[
\frac{\|E_{\psi,c,k}\|_2^2}{N_xN_y}
+\frac{\|E_{U,c,k}\|_2^2}{N_xN_y}
\right].}
\]
Training-time counts must distinguish
\[
\#\text{trajectories}\ne\#\text{operator pairs}\ne\#\text{pixels/grid points}.
\]
Large numbers of target times do not create new independent physical cases.
\end{frame}

\setmodule{Module 4}{4 of 4}
\begin{frame}[t]{Module 4: evaluation and limitation}
\vspace{-0.15em}
Full-field metrics use inverse-transformed predictions:
\[
\epsilon^{(q)}_{L_2}(c,k)=\frac{\|\widehat q_{c,k}-q_{c,k}\|_2}{\|q_{c,k}\|_2+\epsilon},
\qquad q\in\{\psi,U,J,\omega\}.
\]
Derived diagnostics are computed after prediction:
\[
\widehat J=-\nabla^2\widehat\psi,
\qquad
\widehat\omega=\nabla^2\widehat U,
\qquad
E_B,E_K,J_{\max},\Phi_{\rm rec},\text{ spectra, O/X points}.
\]
Evaluation-only residuals may be reported, but not trained:
\[
\Rpsi,\quad \Romega\quad\text{are diagnostics in Module 4.}
\]
Limitation: because $U_0=A_0(\cos x-\cos y)$, the current dataset tests operator learning over transport coefficients, time, and amplitude, not arbitrary initial-state shape generalization.
\end{frame}

% --------------------------------------------------------------------------
% Module 5: four slides
% --------------------------------------------------------------------------

\setmodule{Module 5}{1 of 4}
\begin{frame}[t]{Module 5: physics-informed neural operator}
\vspace{-0.15em}
Module 5 keeps the same full-field FNO map as Module 4:
\[
\mathcal G_\theta:(\psi_0,U_0,\eta,\nu,t)\mapsto(\widehat\psi_t,\widehat U_t).
\]
The model becomes physics-informed when reduced-MHD residuals affect training. The derivative strategy is
\[
\boxed{\text{AD in scalar time }t\quad+\quad\text{Fourier spectral differentiation in }(x,y).}
\]
Predicted derived fields are
\[
\widehat J=-\nabla^2\widehat\psi,
\qquad
\widehat\omega=\nabla^2\widehat U.
\]
This is a physics-informed neural operator/PINO formulation, not the old coordinate-MLP PINN.
\end{frame}

\setmodule{Module 5}{2 of 4}
\begin{frame}[t]{Module 5: full-field residual equations}
\vspace{-0.15em}
For each physics context $(\psi_0,U_0,\eta,\nu,t)$, form complete residual fields:
\[
\boxed{\Rpsi=\widehat\psi_t+\bracket{\widehat U}{\widehat\psi}-\eta\nabla^2\widehat\psi,}
\]
\[
\boxed{\Romega=\widehat\omega_t+\bracket{\widehat U}{\widehat\omega}
-\bracket{\widehat J}{\widehat\psi}-\nu\nabla^2\widehat\omega.}
\]
Spatial operations use the same periodic spectral conventions as the Module 2 solver:
\[
\bracket{f}{g}=f_xg_y-f_yg_x,
\qquad
\widehat{\mathcal B(f,g)}=\mathcal P_{2/3}\mathcal F(f_xg_y-f_yg_x).
\]
The vorticity residual still contains $\nabla^4\widehat U$ mathematically, but it is computed by spectral multiplication rather than four nested coordinate-AD passes.
\end{frame}

\setmodule{Module 5}{3 of 4}
\begin{frame}[t]{Module 5: objective and constraints}
\vspace{-0.15em}
The primary loss is
\[
\boxed{\mathcal L_{\rm total}=w_d\mathcal L_{\rm data}
+w_f\mathcal L_{\rm PDE}+w_0\mathcal L_{\rm IC}+w_g\mathcal L_{\rm gauge},
\qquad \mathcal L_{\rm BC}=0.}
\]
The PDE term is an equation-balanced field norm:
\[
\mathcal L_{\rm PDE}=\lambda_\psi\frac{\|\Rpsi/S_\psi\|_2^2}{N_xN_y}
+\lambda_\omega\frac{\|\Romega/S_\omega\|_2^2}{N_xN_y}.
\]
The operator identity at the initial time may be imposed by
\[
\mathcal L_{\rm IC}=\|\widehat\psi_{t=0}-\psi_0\|^2+\|\widehat U_{t=0}-U_0\|^2.
\]
The stream-function gauge remains
\[
\mathcal L_{\rm gauge}\propto \avg{\widehat U}^{\;2}.
\]
Energy-balance losses are controlled ablations, not silently added to the base model.
\end{frame}

\setmodule{Module 5}{4 of 4}
\begin{frame}[t]{Module 5: controlled evidence that physics helped}
\vspace{-0.15em}
The comparison with Module 4 must hold fixed
\[
\text{trajectories},\quad \text{target times},\quad \text{preprocessing},\quad
\text{FNO capacity},\quad \mathcal L_{\rm data},\quad \text{test cases},\quad \text{metrics}.
\]
Physics helps only if it improves declared outcomes:
\[
\epsilon_{\psi,U},\quad \epsilon_{J,\omega},\quad \|\Rpsi\|,
\quad \|\Romega\|,
\quad E\text{ balance},\quad J_{\max},\quad \Phi_{\rm rec},\quad \text{event timing},\quad \text{data efficiency}.
\]
Failure modes remain plausible:
\[
\begin{array}{lll}
\text{loss imbalance}, & \text{gradient conflict}, & \text{oversmoothing sheets},\\
\text{aliasing mismatch}, & \text{test-case leakage}, & \text{false extrapolation confidence}.
\end{array}
\]
Scientific boundary: the result is still a surrogate of verified synthetic two-field reduced-MHD data. M3D-C1 fidelity requires genuine M3D-C1 arrays and renewed validation.
\end{frame}

\end{document}
